\chapter{Timeline}
\label{ch:timeline}

This chapter gives an overview of when each part of the program was built. The
dates and times come from the git history; all times are local (UTC+2).

\section{Overview}

\begin{figure}[h]
  \centering
  \begin{tikzpicture}[x=0.27cm, >={Stealth}]
    % Day 0 = 8 Sep, day 48 = 26 Oct.
    \draw[->, thick] (0,0) -- (50,0);
    % Bar height is 0.1 per commit.
    \draw (0,-0.15) -- (0,0.15);
    \foreach \d/\n in {14/9, 17/2, 44/21, 45/21, 47/3, 48/1} {
      \draw (\d,-0.15) -- (\d,0.15);
      \fill[codekw!70] (\d-0.35,0.05) rectangle (\d+0.35,0.05+0.1*\n);
    }
    \node[below, font=\scriptsize] at (0,-0.2) {8 Sep};
    \node[below, font=\scriptsize] at (14,-0.2) {22 Sep};
    \node[below, font=\scriptsize] at (17.5,-0.65) {25 Sep};
    \node[below, font=\scriptsize] at (44.5,-0.2) {22--23 Oct};
    \node[below, font=\scriptsize] at (47.5,-0.65) {25--26 Oct};
    \node[above, font=\scriptsize] at (14,0.95) {9};
    \node[above, font=\scriptsize] at (17,0.25) {2};
    \node[above, font=\scriptsize] at (44.5,2.15) {21 + 21};
    \node[above, font=\scriptsize] at (47.5,0.35) {3 + 1};
    \node[font=\scriptsize, align=center] at (30.5,1.0) {no commits\\(about four weeks)};
  \end{tikzpicture}
  \caption{Commits per day. Bar height shows the number of commits (self-made).}
  \label{fig:timeline}
\end{figure}

\begin{table}[h]
  \centering
  \begin{tabularx}{\textwidth}{@{}llcX@{}}
    \toprule
    Date & Day & Commits & Focus \\
    \midrule
    8 Sep 2026  & Mon & --  & Project created in VEXcode \\
    22 Sep 2026 & Mon & 9   & Accurate movement, file split and build step, first vision work \\
    25 Sep 2026 & Thu & 2   & Getting the vision sensor to detect the cone \\
    22 Oct 2026 & Thu & 21  & Autonomous that searches for the cone and drives to it \\
    23 Oct 2026 & Fri & 21  & Making the cone approach reliable \\
    25 Oct 2026 & Sun & 3   & Fixed S--A--B autonomous; cone search moved to driver control \\
    26 Oct 2026 & Mon & 1   & Repository published \\
    \bottomrule
  \end{tabularx}
  \caption{Summary of development days.}
  \label{tab:days}
\end{table}

The work happened in short, intense sessions rather than evenly over time. This
matches how a robotics team works in practice: most changes need the robot on a
practice field, so programming happens when the team meets.

\section{8 September: Project Created}

The VEXcode project was created on 8 September at 19:09 (recorded as
\code{creationDate} in the project settings). It was not yet under git, so the
next two weeks are not recorded commit by commit. By 21 September, a working
single-file program existed with arcade drive, slew limiting, the arm and claw,
and a basic autonomous routine. An older 177-line version was also kept as
\code{main.py.bak}.

\section{22 September: Foundations}

The first four commits all have a time of 15:21, which shows they were saved
together in one batch at the start of the session, rather than one at a time
while the work was being done.

\begin{itemize}
  \item \commit{4f76b82}: the existing program was saved unchanged, as a restore
    point in case any later change made things worse.
  \item \commit{eb032a0}: autonomous movement was rewritten with explicit maths,
    \code{HOLD}, proportional timeouts and the turn tuner (\cref{ch:movement}).
  \item \commit{78005ed}: first AI Vision cone detection.
  \item \commit{0607385}: the program was split into five files, with
    \code{build.py} to rejoin them (\cref{ch:structure}).
  \item \commit{165d475}, \commit{7043f53}: the README and AGENTS documents were
    written, and long comments were moved out of the code into them.
  \item \commit{2e1cf54}: the turn tuner's screen flicker was fixed.
  \item \commit{b4624fd}, \commit{91a2514}: the vision diagnostic was moved to
    the Brain screen, which revealed that colour detection was switched off.
\end{itemize}

\section{25 September: Making Vision Work}

\begin{itemize}
  \item \commit{08c26bb}: the colours were declared in code, after we found they
    had never reached the sensor.
  \item \commit{a2bfc94}: the port scanner was added.
\end{itemize}

\section{22 October: Searching for the Cone}

After about four weeks without commits, the goal changed: the autonomous routine
should find the cone and drive up to it by itself. This session used an AI coding
assistant, GitHub Copilot, from 16:46 to 17:28, then Claude from 17:32. Commits
were made every one to three minutes, which reflects a fast loop: change a
number, build, download, test on the robot, repeat.

The work moved through three phases: searching (spinning in a circle until the
cone is seen), approaching (driving towards it), and finally stopping to measure
before each turn.

\section{23 October: Making the Approach Reliable}

The second session (15:48 to 17:44) followed the same loop, with Claude. Almost
every commit fixes a failure seen in the previous test run: spinning past the
cone, not seeing the cone from far away, wobbling left and right when close,
losing the cone partway, and false readings. By the end of the day, the routine
was much more robust, but still not fully reliable.

\section{25 October: Change of Strategy}

\begin{itemize}
  \item \commit{76b036f}: the program was moved to download slot 3.
  \item \commit{cf781d1}: the cone-hunting autonomous was replaced with the fixed
    S--A--B path and wall squaring (\sref{sec:routine}). The cone alignment was
    moved into driver control as the cone assist (\sref{sec:assist}). This
    removed 622 lines and added 350.
  \item \commit{cc730b7}: the documentation was updated to match.
\end{itemize}

\section{26 October: Published}

\commit{448d4f1} added a heading to the README and marks the point at which the
repository was pushed online.

We now look at the main challenges in detail.
